\documentclass[12pt]{article}
\usepackage[utf8]{inputenc}
\usepackage[T1]{fontenc}
\usepackage{hyperref}
\title{Yambes Language Resources for K-12 Curriculum Development}
\author{Compiled by Claude Sonnet 5.5 (claude-sonnet-5-5) for Andr\'as Kornai}
\date{}
\begin{document}
\maketitle

\section{Introduction}
Yambes (ISO 639-3: ymb; Glottocode yamb1254) is a Torricelli language of East Sepik Province, Papua New Guinea. It is spoken in Yambes village, Dreikikier Rural LLG. Wikipedia gives about 1,100 native speakers (2003). Joshua Project reports that adults still use the language while younger people are shifting to more widely spoken languages. Joshua Project lists the dialects East Yambes, Labuain and West Yambes. Classification is unsettled: Ross (2005) left it unclassified within Torricelli because of sparse data, and Dryer (2024, unpublished) places it in Kombioic.

The searches found only general reference and demographic pages. They found no primers, textbooks, dictionaries, grammars or school programmes for Yambes. A K-12 curriculum would therefore need new documentation work with the community.

\section{Primary Educational Resources}
No primers, textbooks or school materials in Yambes turned up in the searches. The language is not known to be taught in schools.

\section{Linguistic and Reference Materials}
\begin{itemize}
\item Wikipedia, Yambes language (overview, classification, speaker count): \url{https://en.wikipedia.org/wiki/Yambes_language}
\item Crossgram entry for Glottocode yamb1254: \url{https://crossgram.clld.org/languages/yamb1254}
\end{itemize}

\section{Religious and Community Resources}
\begin{itemize}
\item Joshua Project, Yambes people group in Papua New Guinea (demographics and language-shift note): \url{https://joshuaproject.net/people_groups/15972/PP}
\item Joshua Project, language page for ymb: \url{https://joshuaproject.net/languages/ymb}
\end{itemize}

\section{Digital and Audio Resources}
No audio, video or corpus resources specific to Yambes turned up in the searches.

\section{Curriculum Development Context}
Little published material exists for Yambes. Because adults are still first-language speakers while children are shifting away, early priorities would be documenting the language with community members: a practical orthography, a word list, and recorded stories. Teaching materials could be built on those. This should be done with the Yambes community and local education authorities in East Sepik Province.

\end{document}